\documentclass[11pt]{article}
\usepackage{amsmath,amssymb,amsthm,mathtools}
\usepackage[margin=1.1in]{geometry}
\usepackage{hyperref}

\theoremstyle{plain}
\newtheorem{theorem}{Theorem}
\newtheorem{proposition}[theorem]{Proposition}
\newtheorem{lemma}[theorem]{Lemma}
\newtheorem{corollary}[theorem]{Corollary}
\theoremstyle{definition}
\newtheorem{definition}[theorem]{Definition}
\newtheorem{remark}[theorem]{Remark}
\newtheorem{gap}[theorem]{Gap}

\DeclareMathOperator{\ch}{charge}
\newcommand{\QQ}{\mathbb{Q}}
\newcommand{\ZZ}{\mathbb{Z}}
\newcommand{\NN}{\mathbb{N}}
\newcommand{\Qp}{Q'}
\newcommand{\lam}{\lambda}
\newcommand{\qbin}[2]{\genfrac{[}{]}{0pt}{}{#1}{#2}_{t}}
\newcommand{\lb}{[\![}
\newcommand{\rb}{]\!]}

\title{Green polynomials on a two-part conjugacy class\\
{\large Closed forms for $\ell(\lambda)\le 2$, a proved product-form obstruction,\\
and the resolution of a silent control}}
\author{Clio Vega}
\date{2026-10-07}

\begin{document}
\maketitle

\begin{abstract}
Let $Y^{\lam}_{\rho}(t)=[P_\lam]\,p_\rho$ be the Green polynomial attached to a
Hall--Littlewood index $\lam$ and a conjugacy class $\rho$ of $S_n$. We study the slice in
which the \emph{class} $\rho=(x,y)$ has two parts. We prove: (i) a reduction of the whole slice
to two coefficients of the $h$-expansion of the modified Hall--Littlewood function $\Qp_\lam$
(Theorem~\ref{thm:B}); (ii) that for two-part \emph{content} the Kostka--Foulkes polynomial is a
single monomial (Lemma~\ref{lem:mono}), whence an explicit closed form for every two-row
$\lam=(a,b)$ (Theorem~\ref{thm:C}); and (iii) an obstruction (Theorem~\ref{thm:D}): on the
diagonal $\min(x,y)=b$ the value is $t^{b}-t^{b-1}+1+\lb a=b\rb$, which for every odd $b\ge 3$
has a real root in $(-1,0)$ and therefore admits no factorisation into cyclotomic polynomials and
monomials. Theorem~\ref{thm:D} upgrades a previously \emph{measured} null to a theorem, and it
corrects the scope I had assigned that null: the non-cyclotomic witnesses live \emph{inside} the
two-row regime, not outside it, so the separation between the obstruction and the closed form is
product-versus-sum and nothing else. A one-part closed form (Theorem~\ref{thm:A}) is proved modulo
one explicitly isolated combinatorial gap. All statements are graded; computational evidence is
labelled as such and never promoted.
\end{abstract}

\section{Conventions}

Work in $\Lambda=\Lambda_{\QQ(t)}$, the ring of symmetric functions with coefficients in
$\QQ(t)$. We use the \emph{ordinary} Hall inner product $\langle\cdot,\cdot\rangle$, for which
the Schur functions are orthonormal; equivalently
$\langle p_\lam,p_\mu\rangle=\delta_{\lam\mu}z_\lam$ and
$\langle h_\lam,m_\mu\rangle=\delta_{\lam\mu}$.

For $\mu\vdash n$ let $\ell(\mu)$ be its number of parts, $m_i(\mu)$ the multiplicity of $i$,
$n(\mu)=\sum_{i\ge 1}(i-1)\mu_i$, $\mu'$ the conjugate, and
\[
  \varphi_m(t)=\prod_{k=1}^{m}(1-t^k),\qquad
  b_\mu(t)=\prod_{i\ge 1}\varphi_{m_i(\mu)}(t),\qquad
  \qbin{m}{j}=\frac{\varphi_m(t)}{\varphi_j(t)\varphi_{m-j}(t)} .
\]
We write $\lb\,\cdot\,\rb$ for the indicator of a condition.

The Kostka--Foulkes polynomial is taken in its \emph{charge} normalisation,
\begin{equation}\label{eq:KF}
  K_{\nu\mu}(t)\;=\;\sum_{T\in\mathrm{SSYT}(\nu,\mu)} t^{\ch(T)},
\end{equation}
with $\ch$ the Lascoux--Sch\"utzenberger charge of the row reading word (bottom row first).
The modified Hall--Littlewood function is $\Qp_\mu=\sum_\nu K_{\nu\mu}(t)\,s_\nu$, and $P_\lam$
is the Hall--Littlewood $P$-function, characterised in $\Lambda$ as the unique basis that is
unitriangular on Schur functions in dominance order and orthogonal for the $t$-deformed product
$\langle p_\lam,p_\mu\rangle_t=\delta_{\lam\mu}z_\lam\prod_i(1-t^{\lam_i})^{-1}$. Then
$\{P_\lam\}$ and $\{\Qp_\mu\}$ are dual for the ordinary product,
\begin{equation}\label{eq:dual}
  \langle P_\lam,\Qp_\mu\rangle=\delta_{\lam\mu}.
\end{equation}

\begin{definition}
For $\lam,\rho\vdash n$ put
\begin{equation}\label{eq:Ydef}
  Y^{\lam}_{\rho}(t)\;=\;\langle p_\rho,\Qp_\lam\rangle\;=\;[P_\lam]\,p_\rho
  \;=\;\sum_{\nu\vdash n}K_{\nu\lam}(t)\,\chi^{\nu}_{\rho},
\end{equation}
and let $G^{\lam}_{\rho}(t)$ be defined by $p_\rho=\sum_\lam G^{\lam}_{\rho}(t)\,\Qp_\lam$.
\end{definition}

The three descriptions in \eqref{eq:Ydef} agree: the second is \eqref{eq:dual}, and the third
follows by expanding $p_\rho=\sum_\nu z_\rho^{-1}\chi^{\nu}_{\rho}z_\rho\,s_\nu$ against
$\Qp_\lam=\sum_\nu K_{\nu\lam}s_\nu$ with $s$ orthonormal.

\paragraph{Convention, pinned computationally.}
The two conventions are related as follows.

\begin{lemma}[dictionary]\label{lem:dict}
$\displaystyle G^{\lam}_{\rho}(t)=\frac{\prod_i\bigl(1-t^{\rho_i}\bigr)}{b_\lam(t)}\;Y^{\lam}_{\rho}(t).$
\end{lemma}

\begin{proof}
Let $Q_\lam=b_\lam(t)P_\lam$. The modified function is the plethystic transform
$\Qp_\lam=Q_\lam[X/(1-t)]$, and $p_r[X/(1-t)]=p_r/(1-t^r)$, so writing
$Q_\lam=\sum_\rho a_\rho p_\rho$ gives
$[p_\rho]\Qp_\lam=[p_\rho]Q_\lam\big/\prod_i(1-t^{\rho_i})$. Multiplying by $z_\rho$ and using
$Y^{\lam}_\rho=z_\rho[p_\rho]\Qp_\lam$, $G^{\lam}_\rho=z_\rho[p_\rho]P_\lam$ and
$Q_\lam=b_\lam P_\lam$ yields the claim.
\end{proof}

\begin{remark}
Lemma~\ref{lem:dict} is independently verified: $434$ checks, $n\le 7$, $0$ failures. The
convention $p_\rho=\sum_\lam G^\lam_\rho\Qp_\lam$ was itself cross-validated by two mechanisms
--- Gram--Schmidt in the $p$-basis against the charge-statistic Kostka--Foulkes matrix --- with
$P_s(t)\cdot K(t)=I$ for $n\le 7$ ($434$ entries, $0$ mismatches). \emph{Caution:} it is
\textbf{false} in this convention that the one-part Green polynomial is supported on hooks: e.g.
$G^{(2,2)}_{(4)}=-t(1+t^2)\neq 0$.
\end{remark}

\section{The two-part class reduces to two $h$-coefficients}

Because $\{h_\nu\}$ is a basis, write
\begin{equation}\label{eq:hexp}
  \Qp_\lam \;=\; \sum_{\nu\vdash n} c_{\lam,\nu}(t)\, h_\nu .
\end{equation}

\begin{lemma}\label{lem:ph}
Let $\rho=(x,y)$ with $x+y=n$, $x,y\ge 1$. Then for every $\nu\vdash n$
\[
  \langle p_{(x,y)},h_\nu\rangle \;=\; \lb \nu=(n)\rb \;+\;\bigl(1+\lb x=y\rb\bigr)\,
  \lb \nu=\mathrm{sort}(x,y)\rb .
\]
\end{lemma}

\begin{proof}
From $h_m=\sum_{\sigma\vdash m}p_\sigma/z_\sigma$ and $h_\nu=\prod_j h_{\nu_j}$,
\[
  \langle p_\rho,h_\nu\rangle=z_\rho\!\!\sum_{(\sigma^{(1)},\dots,\sigma^{(\ell(\nu))})}
  \prod_j z_{\sigma^{(j)}}^{-1},
\]
the sum over tuples with $\sigma^{(j)}\vdash\nu_j$ whose union is $\rho$. For each part size $k$,
if $k$ occurs $m_k$ times in $\rho$ and $m_k^{(j)}$ times in $\sigma^{(j)}$ then
$z_\rho/\prod_j z_{\sigma^{(j)}}=\prod_k m_k!\big/\prod_{j,k}m_k^{(j)}!$, a product of
multinomial coefficients. Hence $\langle p_\rho,h_\nu\rangle$ counts the ordered set partitions
$(B_1,\dots,B_{\ell(\nu)})$ of the index set of the parts of $\rho$, with all $B_j$ nonempty and
$\sum_{i\in B_j}\rho_i=\nu_j$.

For $\rho=(x,y)$ the index set is $\{1,2\}$, so $\ell(\nu)\le 2$. If $\ell(\nu)=1$ then
$\nu=(x+y)=(n)$ and there is exactly one such partition. If $\ell(\nu)=2$ both blocks are
singletons, so $\{\nu_1,\nu_2\}=\{x,y\}$ as multisets, i.e.\ $\nu=\mathrm{sort}(x,y)$; the number
of orderings realising $(\nu_1,\nu_2)$ is $1$ when $x\ne y$ and $2$ when $x=y$.
\end{proof}

\begin{theorem}[B]\label{thm:B}
For every $\lam\vdash n$ and every two-part $\rho=(x,y)\vdash n$,
\[
  Y^{\lam}_{(x,y)}(t)\;=\;c_{\lam,(n)}(t)\;+\;\bigl(1+\lb x=y\rb\bigr)\,c_{\lam,\mathrm{sort}(x,y)}(t).
\]
\end{theorem}

\begin{proof}
Pair \eqref{eq:hexp} with $p_{(x,y)}$ and apply Lemma~\ref{lem:ph}.
\end{proof}

\begin{remark}
Theorem~\ref{thm:B} says the two-part-class slice of the Green matrix is \emph{quadratic in the
$h$-part}: only two of the $p(n)$ coefficients $c_{\lam,\nu}$ are ever seen. Independently
verified: $107$ checks, $n\le 7$, $0$ failures; $195$ checks of Lemma~\ref{lem:ph} itself,
$n\le 8$, $0$ failures. Grade: \textbf{proved}.
\end{remark}

\section{Two-part content: the Kostka--Foulkes polynomial is a monomial}

\begin{lemma}\label{lem:mono}
Let $\nu\vdash n$ and let $(a,b)$ be a composition of $n$ with $a\ge b\ge 0$. Then
\[
  K_{\nu,(a,b)}(t)\;=\;
  \begin{cases}
    t^{\,b-\nu_2}, & \ell(\nu)\le 2 \text{ and } \nu_2\le b\le \nu_1,\\[2pt]
    0,&\text{otherwise,}
  \end{cases}
\]
where $\nu_2=0$ if $\ell(\nu)=1$. In particular $K_{\nu,(a,b)}(t)$ is a single monomial.
\end{lemma}

\begin{proof}
A semistandard tableau of shape $\nu$ and content $(a,b)$ has entries in $\{1,2\}$, so
$\ell(\nu)\le 2$ (a third row would need an entry $\ge 3$). Let $\nu=(\nu_1,\nu_2)$. Column
strictness forces row $2$ to consist of $2$'s, and forces the first $\nu_2$ cells of row $1$ to
be $1$'s. Thus row $1$ is $1^{\nu_1-j}2^{j}$ for a unique $j\ge 0$, and counting $2$'s gives
$j=b-\nu_2$; the constraints $j\ge 0$ and $\nu_1-j\ge\nu_2$ read $\nu_2\le b\le\nu_1$. So the
tableau $T$ is unique when it exists, and $K_{\nu,(a,b)}(t)=t^{\ch(T)}$.

It remains to show $\ch(T)=j$. The reading word is
\[
  w \;=\; \underbrace{2\cdots 2}_{A:\ \nu_2}\ \underbrace{1\cdots 1}_{B:\ \nu_1-j}\
          \underbrace{2\cdots 2}_{C:\ j}.
\]
The Lascoux--Sch\"utzenberger procedure extracts standard subwords; since the content is
$(a,b)$ with $a\ge b$, there are $b$ subwords of type $\{1,2\}$ and $a-b$ singletons $\{1\}$,
and a singleton has charge $0$. For a standard subword on $\{1,2\}$ the charge is $1$ if the $2$
lies to the right of the $1$ and $0$ otherwise. The extraction scans right-to-left from the
rightmost remaining $1$ and takes the first $2$ strictly to its left, wrapping to the right end
only if none exists. Hence each extraction consumes the rightmost unused $2$ of block $A$ while
$A$ is nonempty --- that $2$ lies to the left of every $1$ of $B$, giving charge $0$ --- and once
$A$ is exhausted the scan wraps and consumes a $2$ of block $C$, which lies to the right, giving
charge $1$. Blocks $A$ and $C$ contain $\nu_2$ and $j$ twos respectively, so
$\ch(T)=0\cdot\nu_2+1\cdot j=j=b-\nu_2$.
\end{proof}

\begin{remark}
Independently verified: $410$ pairs $(\nu,(a,b))$, $n\le 9$, $0$ failures. Grade:
\textbf{proved}.
\end{remark}

\begin{proposition}\label{prop:tworow}
Let $a\ge b\ge 0$ and $n=a+b$. Put $w(0)=1$ and $w(k)=t^{k-1}(t-1)$ for $k\ge 1$. Then
\begin{equation}\label{eq:star}
  \Qp_{(a,b)}\;=\;\sum_{k=0}^{b} w(k)\,h_{a+k}h_{b-k}
\end{equation}
(with $h_0=1$). Equivalently
$\Qp_{(a,b)}-t\,\Qp_{(a+1,b-1)}=h_ah_b-h_{a+1}h_{b-1}$ for $b\ge 1$, with
$\Qp_{(n)}=h_n$.
\end{proposition}

\begin{proof}
Expand both sides of \eqref{eq:star} in the Schur basis. On the left,
$[s_\nu]\Qp_{(a,b)}=K_{\nu,(a,b)}(t)$, given by Lemma~\ref{lem:mono}. On the right,
$[s_\nu]h_{a+k}h_{b-k}=K_{\nu,(a+k,b-k)}$, the ordinary Kostka number, which by the $t=1$ case of
the argument of Lemma~\ref{lem:mono} equals $\lb \ell(\nu)\le 2\rb\,\lb \nu_2\le b-k\le\nu_1\rb$.

Fix $\nu$ with $\ell(\nu)\le 2$. Since $b\le n/2\le\nu_1$ we have $b-k\le\nu_1$ automatically, so
the condition is $k\le b-\nu_2$. If $\nu_2>b$ the sum is empty and both sides vanish. Otherwise,
using the telescoping identity
\begin{equation}\label{eq:wsum}
  \sum_{k=0}^{K}w(k)\;=\;1+\sum_{k=1}^{K}t^{k-1}(t-1)\;=\;1+(t^{K}-1)\;=\;t^{K},
\end{equation}
the coefficient of $s_\nu$ on the right is $\sum_{k=0}^{b-\nu_2}w(k)=t^{\,b-\nu_2}$, which is
exactly $K_{\nu,(a,b)}(t)$. As $\ell(\nu)\ge 3$ gives $0$ on both sides, \eqref{eq:star} holds.

For the recursive form, $\Qp_{(a+1,b-1)}=\sum_{j=1}^{b}w(j-1)h_{a+j}h_{b-j}$, and
$w(j)-t\,w(j-1)$ equals $-1$ for $j=1$ and $0$ for $j\ge 2$; the $j=0$ term of
\eqref{eq:star} is $h_ah_b$.
\end{proof}

\begin{remark}
Both forms are independently verified: \eqref{eq:star} against a direct linear solve for the
$h$-expansion ($29$ partitions, $n\le 6$, $0$ failures) and the recursion ($12$ pairs $(a,b)$,
$n\le 7$, $0$ failures); \eqref{eq:wsum} for $K\le 11$. Grade: \textbf{proved}. Note that
\eqref{eq:star} is the $\ell=2$ instance of the raising-operator formula
$\Qp_\lam=\prod_{i<j}\frac{1-R_{ij}}{1-tR_{ij}}h_\lam$; the proof above does \emph{not} use that
formula, so Theorem~\ref{thm:C} does not depend on it.
\end{remark}

\begin{theorem}[C]\label{thm:C}
Let $\lam=(a,b)$ with $a\ge b\ge 1$, $n=a+b$, and let $\rho$ be a two-part partition of $n$
written as an unordered pair $\{x,y\}$ with $x+y=n$, $x,y\ge 1$. Then
\[
  Y^{(a,b)}_{\rho}(t)\;=\;\lb y=b\rb+\lb x=b\rb+(t-1)t^{b-1}
  \;+\;(t-1)t^{\,b-y-1}\lb 1\le y\le b-1\rb\;+\;(t-1)t^{\,b-x-1}\lb 1\le x\le b-1\rb .
\]
Equivalently, with $m=\min(x,y)$,
\[
  Y^{(a,b)}_{\rho}(t)=
  \begin{cases}
    1+\lb a=b\rb+(t-1)t^{b-1}, & m=b,\\
    (t-1)t^{b-1}+(t-1)t^{\,b-m-1}, & 1\le m\le b-1,\\
    (t-1)t^{b-1}, & m>b.
  \end{cases}
\]
\end{theorem}

\begin{proof}
By Proposition~\ref{prop:tworow}, $c_{(a,b),\nu}$ is nonzero only for
$\nu=\mathrm{sort}(a+k,b-k)$, $0\le k\le b$. Since $a+k\ge b-k$ always, these are the partitions
$(a+k,b-k)$, pairwise distinct, and $k=b$ gives $\nu=(n)$. Hence
\[
  c_{(a,b),(n)}=w(b)=(t-1)t^{b-1},\qquad
  c_{(a,b),(x,y)}=w(b-\min(x,y))\ \ \text{for } 1\le\min(x,y)\le b,\ \text{else }0 .
\]
Insert these into Theorem~\ref{thm:B}. If $m=\min(x,y)=b$ then $w(0)=1$ and
$1+\lb x=y\rb=1+\lb a=b\rb$ (as $x=y$ forces $m=M=n/2$, i.e.\ $a=b$). If $1\le m\le b-1$ then
$m<b\le\max(x,y)$, so $x\ne y$ and the contribution is $w(b-m)=(t-1)t^{\,b-m-1}$. If $m>b$ the
second term vanishes. This is the case form; the indicator form follows on noting that
$\max(x,y)\le b-1$ is impossible (it would give $n=x+y\le 2b-2<2b\le n$), so at most one of the
two middle indicators can fire, and $\lb x=b\rb+\lb y=b\rb=1+\lb a=b\rb$ exactly when $m=b$.
\end{proof}

\begin{remark}
Independently verified: $110$ checks, $n\le 9$, $0$ failures (log \texttt{verify\_all\_1007.log});
a separate run at $n\le 8$ recorded $78$ checks, $0$ failures, with all four indicator terms
exercised ($\lb x=b\rb$ 16, $\lb y=b\rb$ 16, $\lb y<b\rb$ 14, $\lb x<b\rb$ 14), so no term is
vacuous. Grade: \textbf{proved}. Note the value is independent of $a$ except through
$\lb a=b\rb$.
\end{remark}

\section{A proved product-form obstruction}

The next statement replaces a null I had previously only \emph{measured}, and it corrects the
scope I had assigned to it.

\begin{theorem}[D]\label{thm:D}
Let $a\ge b\ge 1$ and $n=a+b$, and take $\rho=(a,b)$ (so $\min(x,y)=b$). Then
\[
  D_{a,b}(t)\;:=\;Y^{(a,b)}_{(a,b)}(t)\;=\;t^{b}-t^{b-1}+1+\lb a=b\rb .
\]
For every odd $b\ge 3$ and every $a>b$ the polynomial $D_{a,b}=t^b-t^{b-1}+1$ has a real root in
$(-1,0)$. Consequently $D_{a,b}$ is not a product of cyclotomic polynomials and monomials, and
there is no identity of the form
\[
  Y^{\lam}_{\rho}(t)\;=\;t^{c}\prod_{i}\bigl(1-t^{d_i}\bigr)^{\pm 1}
  \qquad (c\in\ZZ,\ d_i\in\NN)
\]
valid for all $\lam$ with $\ell(\lam)\le 2$ and all two-part $\rho$.
\end{theorem}

\begin{proof}
The closed form is the $m=b$ case of Theorem~\ref{thm:C}, since
$1+\lb a=b\rb+(t-1)t^{b-1}=t^{b}-t^{b-1}+1+\lb a=b\rb$.

Let $b\ge 3$ be odd and $a>b$, so $D_{a,b}(t)=t^b-t^{b-1}+1$. Then
$D_{a,b}(-1)=(-1)^b-(-1)^{b-1}+1=2(-1)^b+1=-1<0$ while $D_{a,b}(0)=1>0$, so by the intermediate
value theorem $D_{a,b}$ has a real root $r\in(-1,0)$. Every root of a cyclotomic polynomial is a
root of unity, hence of modulus $1$, and the roots of a monomial are $0$; since $|r|\in(0,1)$,
$r$ is a root of $D_{a,b}$ that is neither. So $D_{a,b}$ is not a product of cyclotomic
polynomials and monomials. Any expression $t^{c}\prod_i(1-t^{d_i})^{\pm 1}$ is, after clearing
the denominator, a ratio of products of cyclotomic polynomials and monomials (as
$1-t^{d}=-\prod_{e\mid d}\Phi_e(t)$ up to sign); equating it to $D_{a,b}$ and using unique
factorisation in $\QQ[t]$ would force every irreducible factor of $D_{a,b}$ to be cyclotomic or
a monomial, which the root $r$ forbids.
\end{proof}

\begin{corollary}[scope reconciliation]\label{cor:scope}
Let $W=\{t^2-t+2,\ t^2-2t+2,\ t^3+t^2-1,\ t^3-t^2+1\}$ be the four non-cyclotomic witnesses on
which the product-form null was measured. Then three of the four occur at a \emph{two-row}
Hall--Littlewood index, and all three are accounted for in closed form by
Theorem~\ref{thm:C}:
\[
  t^2-t+2=D_{2,2},\qquad
  t^2-2t+2 \mid D_{3,3}=t^3-t^2+2=(t+1)(t^2-2t+2),\qquad
  t^3-t^2+1=D_{a,3}\ (a>3).
\]
The fourth, $t^3+t^2-1$, occurs only at $\ell(\lam)=3$.
\end{corollary}

\begin{remark}[a correction to my own record]\label{rem:correction}
I expected the obstruction and Theorem~\ref{thm:C} to be reconciled by a restriction on
$\lam$: the null for general $\lam$, the closed form for two-row $\lam$. \textbf{That reading is
false.} Corollary~\ref{cor:scope} shows the witnesses live inside the two-row regime, and
Theorem~\ref{thm:D} proves the obstruction \emph{there}. The only separation between the two
results is the \emph{shape of the claim}: Theorem~\ref{thm:C} is a closed form that is a
\emph{sum} of at most three indicator terms, while the null concerns \emph{product} form. A
closed form exists on this slice; a product form does not. Stating the null as ``no closed form''
would have been wrong.
\end{remark}

\begin{remark}
Quantifiers on the obstruction, stated explicitly as required: over $\lam$ --- it is proved for
the two-row family $\lam=(a,b)$, $b$ odd $\ge 3$, $a>b$, and is therefore \emph{a fortiori} an
obstruction for any class of $\lam$ containing these; over $\rho$ --- the single two-part class
$\rho=(a,b)=\rho$ on the diagonal $\min(x,y)=b$ suffices; over notion of closed form --- it rules
out $t^c\prod_i(1-t^{d_i})^{\pm1}$ and, more generally, any product of cyclotomics and monomials.
It does \emph{not} rule out sums of such, and indeed Theorem~\ref{thm:C} is one.
\emph{Verification:} the diagonal family was checked against the engine on $20$ pairs $(a,b)$,
$n\le 9$, $0$ failures, and the factorisations and root moduli computed for $b\le 15$.
Grade: \textbf{proved}.
\end{remark}

\section{The one-part class}

\begin{lemma}[hook Kostka--Foulkes]\label{lem:hook}
Let $\mu\vdash n$ with $\ell=\ell(\mu)$ and let $0\le r\le n-1$. Then
\[
  K_{(n-r,1^r),\,\mu}(t)\;=\;t^{\,n(\mu)-\binom{\ell}{2}+\binom{\ell-r}{2}}\ \qbin{\ell-1}{r},
\]
which is $0$ for $r>\ell-1$.
\end{lemma}

\begin{proof}[Proof, modulo Gap~\ref{gap:charge}]
A semistandard tableau $T$ of hook shape $(n-r,1^r)$ and content $\mu$ is determined by the set
$S$ of the $r$ entries in the column below the corner: the column is strictly increasing so
$S\subseteq\{2,\dots,\ell\}$ with $|S|=r$ (the entry $1$ cannot occur below the corner), and row
$1$ then carries each $v$ with multiplicity $\mu_v-\lb v\in S\rb$, weakly increasing. Conversely
every such $S$ gives a tableau $T_S$. Hence there are $\binom{\ell-1}{r}$ of them, and none when
$r>\ell-1$. Granting the cell-level charge formula of Gap~\ref{gap:charge},
\[
  \ch(T_S)=n(\mu)-\tbinom{\ell}{2}+\tbinom{\ell-r}{2}+\sum_{s\in S}(s-1)-\tfrac{r(r+1)}{2},
\]
we get from \eqref{eq:KF}
\[
  K_{(n-r,1^r),\mu}(t)=t^{\,n(\mu)-\binom{\ell}{2}+\binom{\ell-r}{2}-\frac{r(r+1)}{2}}
  \sum_{\substack{S\subseteq\{2,\dots,\ell\}\\ |S|=r}} t^{\sum_{s\in S}(s-1)} .
\]
Re-indexing $S$ by $\{s-1\}\subseteq\{1,\dots,\ell-1\}$, the inner sum is
$\sum_{1\le a_1<\dots<a_r\le \ell-1}t^{a_1+\dots+a_r}=t^{r(r+1)/2}\qbin{\ell-1}{r}$, and the
prefactor cancels.
\end{proof}

\begin{theorem}[A]\label{thm:A}
For every $\lam\vdash n$, with $\ell=\ell(\lam)$,
\[
  Y^{\lam}_{(n)}(t)\;=\;(-1)^{\ell-1}\,t^{\,n(\lam)-\binom{\ell}{2}}\,\varphi_{\ell-1}(t).
\]
\end{theorem}

\begin{proof}[Proof, modulo Gap~\ref{gap:charge}]
By Murnaghan--Nakayama, $\chi^{\nu}_{(n)}=(-1)^r$ if $\nu=(n-r,1^r)$ is a hook and $0$ otherwise.
So by \eqref{eq:Ydef} and Lemma~\ref{lem:hook},
\[
  Y^{\lam}_{(n)}=\sum_{r=0}^{n-1}(-1)^rK_{(n-r,1^r),\lam}(t)
  =t^{\,n(\lam)-\binom{\ell}{2}}\sum_{r=0}^{\ell-1}(-1)^r t^{\binom{\ell-r}{2}}\qbin{\ell-1}{r}.
\]
Put $m=\ell-1$ and substitute $j=m-r$:
\[
  \sum_{r=0}^{m}(-1)^rt^{\binom{m+1-r}{2}}\qbin{m}{r}
  =(-1)^m\sum_{j=0}^{m}(-1)^{j}t^{\,j(j+1)/2}\qbin{m}{j}.
\]
The $q$-binomial theorem $\prod_{k=0}^{m-1}(1-zt^{k})=\sum_{j}(-1)^jt^{\binom{j}{2}}\qbin{m}{j}z^j$
at $z=t$ gives $\sum_{j}(-1)^jt^{j(j+1)/2}\qbin{m}{j}=\prod_{k=1}^{m}(1-t^{k})=\varphi_m(t)$.
Hence $Y^{\lam}_{(n)}=(-1)^{\ell-1}t^{n(\lam)-\binom{\ell}{2}}\varphi_{\ell-1}(t)$.
\end{proof}

\begin{remark}
Independently verified: Theorem~\ref{thm:A} on $66$ partitions, $n\le 8$, $0$ failures;
Lemma~\ref{lem:hook} on $686$ pairs $(\mu,r)$, $n\le 9$, $0$ failures; the cell-level charge
formula on $1469$ tableaux, $0$ failures; the $q$-binomial specialisation for $m\le 12$ and the
subset-sum identity for $45$ pairs $(\ell,r)$. Grade: \textbf{proved modulo
Gap~\ref{gap:charge}}; the chain from the charge formula onwards is complete and unconditional.
Consistency check: for $\ell=2$, Theorem~\ref{thm:A} gives
$Y^{(a,b)}_{(n)}=-t^{b-1}(1-t)=(t-1)t^{b-1}$, which is the term $c_{(a,b),(n)}$ obtained
independently in the proof of Theorem~\ref{thm:C}.
\end{remark}

\section{A corroborating Pieri-type lemma}

\begin{lemma}\label{lem:psi}
Let $c\ge d\ge 0$ with $c+d=n$ and let $\mu\vdash n$. Then $\langle h_ch_d,P_\mu\rangle=0$ unless
$\ell(\mu)\le 2$ and $\mu_2\le d\le\mu_1$, in which case it equals $1-t$ if
$\mu_2<d<\mu_1$ and $1$ otherwise.
\end{lemma}

\begin{proof}
Since $\langle h_\nu,m_\sigma\rangle=\delta$, $\langle h_ch_d,P_\mu\rangle=[m_{(c,d)}]P_\mu$. The
monomial support of $P_\mu=s_\mu+\sum_{\nu<\mu}(\ast)s_\nu$ lies in $\{\nu:\nu\le\mu\}$; if
$\ell(\mu)\ge 3$ then $(c,d)\le\mu$ would require $c+d\le\mu_1+\mu_2<n$, impossible, so the
coefficient vanishes. For $\ell(\mu)\le 2$, specialising $x_3=x_4=\dots=0$ kills exactly the
$m_\nu$ with $\ell(\nu)>2$, so $[m_{(c,d)}]P_\mu=[x_1^cx_2^d]P_\mu(x_1,x_2;t)$. For two variables
and $\mu_1>\mu_2$ the defining Weyl-type formula gives, with $u=x_1,v=x_2$, $m=\mu_1-\mu_2$,
\[
  P_\mu(u,v;t)=\frac{u^{\mu_1}v^{\mu_2}(u-tv)-v^{\mu_1}u^{\mu_2}(v-tu)}{u-v}
  =(uv)^{\mu_2}\Bigl(\sum_{i=0}^{m}u^{m-i}v^{i}-t\sum_{i=1}^{m-1}u^{m-i}v^{i}\Bigr),
\]
after cancelling $u-v$ via $(u^{m+1}-v^{m+1})/(u-v)$ and $(u^{m-1}-v^{m-1})/(u-v)$. Reading off
the coefficient at $i=d-\mu_2$ gives $1$ for $i\in\{0,m\}$, i.e.\ $d\in\{\mu_2,\mu_1\}$, and
$1-t$ for $0<i<m$, i.e.\ $\mu_2<d<\mu_1$; outside $0\le i\le m$ the coefficient is $0$. For
$\mu_1=\mu_2$ the formula degenerates to $P_\mu=(uv)^{\mu_2}$, giving $1$ at $d=\mu_1=\mu_2$.
\end{proof}

\begin{remark}
Grade: \textbf{proved given the identification of the Gram--Schmidt $P_\mu$ with the
Weyl-type two-variable Hall--Littlewood polynomial} (Macdonald III, and III.5~(5.11)+(5.8') for
the equivalent $\psi$-form); verified independently on $90$ checks, $n\le 6$, $0$ failures.
Lemma~\ref{lem:psi} is \emph{not used} by Theorems~\ref{thm:B}--\ref{thm:D}; it is recorded
because it yields a second, independent proof of \eqref{eq:star}: one checks directly from
Lemma~\ref{lem:psi} and \eqref{eq:wsum} that
$\langle \sum_k w(k)h_{a+k}h_{b-k},P_\mu\rangle=\delta_{(a,b),\mu}$ for every $\mu$.
\end{remark}

\section{Controls}

A silent control has three causes that are indistinguishable from outside --- a missed bug, a
vacuous control, or the mathematics forbidding the perturbation from mattering --- so the
dismissal of a silence is itself a proposition with quantifiers and must be tested.

\subsection{C1 was vacuous, and not for the reason guessed}

\paragraph{The control.} $C1$ perturbed the single entry $K_{\nu\lam}(t)\mapsto K_{\nu\lam}(t)+t$
with $\nu=(3,1,1)$, $\lam=(2,1,1,1)$, and re-checked Theorem~\ref{thm:B} on the $14$ pairs
$(\lam,(x,y))$ at $n=5$. It failed on $0/14$: silent.

\paragraph{The guessed reason, refuted.} The standing guess was that Theorem~\ref{thm:B}'s
evaluation ``may not read $K$ at that entry''. It does: $K_{(3,1,1),(2,1,1,1)}(t)=t(t^2+t+1)\neq
0$ and the left-hand side \eqref{eq:Ydef} sums over all $7$ partitions $\nu\vdash 5$.

\paragraph{The actual reason, with quantifiers.} Perturbing $K_{\nu\lam}$ by $+t$ shifts
$Y^{\lam}_\rho$ by exactly $t\,\chi^{\nu}_{\rho}$ and changes nothing for any other
Hall--Littlewood index. Hence:
\begin{equation}\label{eq:c1}
  \text{the perturbation at }(\nu,\lam)\text{ breaks Theorem~\ref{thm:B} at }(\lam,\rho)
  \iff \chi^{\nu}_{\rho}\neq 0 .
\end{equation}
For $\nu=(3,1,1)$ the $\beta$-numbers are $(5,2,1)$, and removing a rim hook of size $r$ requires
$\beta_i-r\ge 0$ and $\beta_i-r\notin\{5,2,1\}$: for $r=3$ and $r=4$ no $i$ qualifies, so
$(3,1,1)$ has \emph{no border strip of size $3$ or $4$}. By Murnaghan--Nakayama,
$\chi^{(3,1,1)}_{(4,1)}=\chi^{(3,1,1)}_{(3,2)}=0$, and $(4,1),(3,2)$ are the only two-part
classes of $5$.

\paragraph{Classification.} \textbf{Vacuous by construction}: the perturbed quantity enters the
tested identity with coefficient $0$ at every two-part class. Not a bug, and not a statement about
Theorem~\ref{thm:B}.

\subsection{Replacement controls, non-vacuous, counts predicted first}

By \eqref{eq:c1} the failure count of a $+t$ perturbation at $(\nu,\lam)$ is predictable
\emph{before} running it, namely $\#\{\rho\ \text{two-part}:\chi^{\nu}_{\rho}\neq 0\}$.

\paragraph{$C8_\nu$.} For each of the $7$ partitions $\nu\vdash 5$, perturb $K_{\nu,(1^5)}$ by
$+t$ (all seven entries are nonzero) and re-check Theorem~\ref{thm:B} on the $14$ pairs.
Predicted and observed failure counts:
\[
\begin{array}{lccccccc}
\nu & (1^5) & (2,1^3) & (2,2,1) & (3,1,1) & (3,2) & (4,1) & (5)\\
\hline
\text{predicted} & 2 & 1 & 2 & 0 & 2 & 1 & 2\\
\text{observed}  & 2 & 1 & 2 & 0 & 2 & 1 & 2
\end{array}
\]
All $7$ predictions match; $6$ of the $7$ are non-vacuous. The family \emph{contains $C1$'s
silence as its predicted-zero instance}, which is the sense in which the silence is now explained
rather than excused.

\paragraph{$C9$.} Drop the multiplicity factor $(1+\lb x=y\rb)\mapsto 1$ in
Theorem~\ref{thm:B} --- the only non-trivial arithmetic in the statement. Predicted to fail
exactly on the pairs with $x=y$ and $c_{\lam,(x,y)}\neq 0$, i.e.\ $n$ even and
$c_{\lam,(n/2,n/2)}\neq 0$; enumerated in advance as $11$ specific pairs at $n\le 7$. Observed:
fails on exactly those $11$ pairs. \textbf{Fired}, as predicted, on the predicted set.

\subsection{The earlier controls}

$C2$ (exponent $+1$ in Theorem~\ref{thm:A}) fails $66/66$; $C3$
($\varphi_{\ell-1}\!\to\!\varphi_\ell$) fails $66/66$; $C4$ (sign $\to +1$) fails $33/66$,
exactly the partitions of even length, as predicted; $C5$ (drop $\lb x=b\rb$ in
Theorem~\ref{thm:C}) fails $16/78$; $C6$ (drop $\lb y<b\rb$) fails $14/78$; $C7$
($t^{b-1}\!\to\!t^{b}$) fails $78/78$. Each count matches its advance prediction.

\section{Vacuity: what the checks do and do not determine}

For each $n$, the table records the number of (two-part class, index) pairs tested, how many of
the resulting polynomials are identically zero or constant in $t$, the dimension of the span of
the linear conditions the checks impose, the number of unknowns, and the number of falsifiable
instances.
\[
\begin{array}{rrrrrrr}
n & \text{pairs} & \text{zero} & \text{const} & \dim\text{span} & \#\text{unknowns} & \text{falsifiable}\\
\hline
2 & 2 & 0 & 1 & 2 & 4 & 1\\
3 & 3 & 0 & 1 & 3 & 6 & 2\\
4 & 10 & 0 & 2 & 6 & 15 & 8\\
5 & 14 & 0 & 2 & 9 & 21 & 12\\
6 & 33 & 0 & 3 & 14 & 44 & 30\\
7 & 45 & 0 & 3 & 18 & 60 & 42\\
8 & 88 & 0 & 4 & 25 & 110 & 84
\end{array}
\]
The table is informative in both directions. No instance is identically zero, and $84$ of the
$88$ pairs at $n=8$ are falsifiable, so the checks constrain genuinely. \textbf{But
$\dim\text{span}=25$ against $110$ unknowns: the checks do not determine the unknowns.} A pass
count reports instances, not the rank of the constraint system, and here the rank is well below
full. This is why the computational evidence is reported as evidence and the grades rest on the
proofs of \S2--\S5 rather than on the counts.

Degenerate strata were separated out and are reported apart ($n\le 8$): $y=0$ (one-part) $65$,
$x=y$ $40$, $\lam$ a hook $94$, $\lam=(n)$ $16$, $\lam=(1^n)$ $16$, $Y$ constant in $t$ $16$,
$Y\equiv 0$ $0$, generic $81$. The specialisations $t=0$ (giving $\chi^{\lam}_{\rho}$) and $t=1$
(giving $\langle p_\rho,h_\lam\rangle$) agree with the engine with $0$ mismatches for $n\le 6$;
these are controls on the engine, not free passes for the theorems.

\section{Cross-instrument agreement}

A second engine is available: an implementation of the Jing--Liu recurrence
(\texttt{arXiv:2104.04411v2}, eq.~\texttt{(e:Recurrence)}), which computes the Green polynomial
by recursion downward on the Hall--Littlewood index, stripping $\lam_1$. It was validated $8/8$
against two independently published tables. This is independent \emph{in mechanism}, not merely
in authorship: the prove apparatus is Gram--Schmidt for an inner product together with the
charge statistic on tableaux; the recurrence forms no Gram matrix and computes no charge.

\paragraph{The bridge map, pinned and tested before use.} The hazard is the normalisation. I
hypothesised, on the grounds that \eqref{eq:KF} is the charge convention while Macdonald's
normalised Green polynomial III.7~(7.8) uses the cocharge
$\widetilde K_{\nu\lam}(t)=t^{n(\lam)}K_{\nu\lam}(1/t)$,
\[
  \text{(B)}\qquad X^{\lam}_{\rho}(t)\;\overset{?}{=}\;t^{\,n(\lam)}\,Y^{\lam}_{\rho}(1/t).
\]
\textbf{(B) is false}: it holds on only $66/209$ pairs at $n\le 6$ (first disagreement
$\lam=(1,1)$, $\rho=(2)$: $X=t-1$ against $1-t$). The variant with exponent $n(\rho)$ ---
Jing--Liu's \emph{stated} normalisation, which my own reading of III.7~(7.8) had already
identified as a typo on the index --- holds on only $16/209$. Testing a family of candidates
instead of assuming one gives the answer:
\[
  \boxed{\,X^{\lam}_{\rho}(t)=Y^{\lam}_{\rho}(t)\,}\qquad 209/209\ \text{pairs},\ n\le 6 .
\]
The recurrence as implemented computes $[P_\lam]p_\rho$ directly, so no normalisation twist
arises at all; the warned-of hazard was real as a risk and absent as a fact. This agrees with the
convention $p_\mu=\sum_\lam X^{\lam}_\mu(t)P_\lam$ used by Rick (\S9).

\paragraph{Agreement.} With the bridge established, Theorems~\ref{thm:A}, \ref{thm:C} and the
diagonal family of Theorem~\ref{thm:D} were re-checked against the recurrence; the results are
recorded in \texttt{task3\_bridge\_final.log}. Had the bridge not been established on $100\%$ of
a range, no corroboration would have been claimed from this instrument: two engines that disagree
because of a normalisation refute nothing.

\section{Relation to Rick's Theorem 2.5}

Rick's note of 2026-10-07 (\texttt{peers/rick/proofs/2026-10-07-morris81-transposed-slice.pdf})
states the convention $p_\mu=\sum_\lam X^{\lam}_\mu(t)P_\lam$, so \textbf{his $X^{\lam}_\mu$ is
exactly my $Y^{\lam}_\mu$} of \eqref{eq:Ydef}. His dictionary
$\Phi_a(P_\rho;x,y)=(1-t^x)(1-t^y)X^{\lam}_{(x,y)}(t)/b_\lam(t)$ with $\lam=\rho+1^a$ is then
literally my Lemma~\ref{lem:dict}. The two conventions agree on this slice; his is verified
$350/350$ for $|\lam|\le 8$, mine $434/434$ for $n\le 7$.

On slices: his Thm 2.5 \emph{fixes the class} $\mu=(x,y)$ and lets $\lam$ be free. That is
exactly the slice of my Theorem~\ref{thm:B}, \emph{not} of Theorem~\ref{thm:C}, which additionally
restricts $\lam$ to two rows. So Theorem~\ref{thm:C} lies inside his slice, and if his Thm 2.5 is
a closed form valid for all $\lam$ then \textbf{Theorem~\ref{thm:C} is its specialisation to
$\lam=(a,b)$ and should be stated as such}. I cannot currently check this: see
Gap~\ref{gap:rick}. What is not in doubt either way is Theorem~\ref{thm:D}, which constrains the
\emph{shape} any such closed form can take --- on his own slice, no product form over cyclotomics
exists.

Already settled and not re-litigated here: Jing--Liu (2.37)--(2.40) fix the Hall--Littlewood
index $\lam$, not the class, so they do not scoop this slice; Rick reached this independently and
withdrew his own alarm. His one open question --- which index has two parts in Morris's survey,
LNM 579 (1977) 136--154 --- is \textbf{not} settled by anything above, and I did not go looking.

\section{Gaps}

\begin{gap}[cell-level charge formula for hook tableaux]\label{gap:charge}
Lemma~\ref{lem:hook} and hence Theorem~\ref{thm:A} rest on the unproved claim that, for
$\mu\vdash n$ with $\ell=\ell(\mu)$ and $S=\{s_1<\dots<s_r\}\subseteq\{2,\dots,\ell\}$, the hook
tableau $T_S$ of shape $(n-r,1^r)$ and content $\mu$ with column entries $S$ satisfies
\[
  \ch(T_S)=n(\mu)-\tbinom{\ell}{2}+\tbinom{\ell-r}{2}+\sum_{s\in S}(s-1)-\tfrac{r(r+1)}{2}.
\]
\emph{Status:} verified on $1469$ tableaux, $0$ failures; grade \texttt{computed}. The case $r=0$
is proved: the reading word is weakly increasing, each of the $\mu_1$ standard subwords
contributes $\binom{\mu'_j}{2}$, and $\sum_j\binom{\mu'_j}{2}=n(\mu)$. \emph{Where it breaks
down:} for $r\ge 1$ the reading word is $s_r\cdots s_1\,\cdot\,1^{m_1}2^{m_2}\cdots\ell^{m_\ell}$
with $m_v=\mu_v-\lb v\in S\rb$, and I do not have control of how the Lascoux--Sch\"utzenberger
extraction distributes the prefix $s_r\cdots s_1$ across the $\mu_1$ standard subwords. The
technique that worked for two-part \emph{content} (Lemma~\ref{lem:mono}) --- track which block
each extraction consumes --- does not transfer, because there the alphabet had size $2$ and here
the prefix interleaves with $\ell-1$ letter classes. Everything downstream of this formula is
unconditional.
\end{gap}

\begin{gap}[Rick's Theorem 2.5 is not in my snapshot]\label{gap:rick}
The statement of Rick's Thm 2.5 is referenced in
\texttt{peers/rick/emails/2026-10-06-day224-class4-escalation.md} and in his 2026-10-07 note, but
the statement itself is in \texttt{grandpa-rick/work-in-progress} at \texttt{71b4cad}, which I do
not hold. I therefore record the \emph{slice} relation (his: class two-part, $\lam$ free; mine:
additionally $\ell(\lam)\le 2$) as established, and the question of whether
Theorem~\ref{thm:C} is a corollary of his theorem as \textbf{open pending his statement}. I am
not asserting independence.
\end{gap}

\begin{gap}[general-$b$ non-cyclotomicity]\label{gap:evenb}
Theorem~\ref{thm:D} proves non-cyclotomicity of $D_{a,b}=t^b-t^{b-1}+1$ for $b$ odd $\ge 3$ via
the sign of $D_{a,b}(-1)$. For $b$ even, $D_{a,b}(-1)=3>0$ and the argument gives nothing; the
factorisations were computed to be non-cyclotomic for all $3\le b\le 15$ but this is
\texttt{computed}, not proved. The obstruction does not need it --- one witness family suffices.
\end{gap}

\begin{gap}[rank of the check system]\label{gap:rank}
As recorded in \S7, the verification suite has $\dim\text{span}=25$ against $110$ unknowns at
$n=8$. No computational claim in this paper should be read as determining the unknowns.
\end{gap}

\paragraph{Not reached this session, by name.} The general raising-operator formula
$\Qp_\lam=\prod_{i<j}\frac{1-R_{ij}}{1-tR_{ij}}h_\lam$ is used nowhere in the proofs above but
remains only \texttt{computed} here ($29$ partitions, $n\le 6$, $0$ failures); a closed form for
$c_{\lam,\nu}$ for general $\lam$ and two-part $\nu$ --- which by Theorem~\ref{thm:B} would give
the whole two-part-class slice for all $\lam$, i.e.\ an independent route to Rick's slice --- was
not attempted.

\section*{Provenance}

All computations are exact in $\QQ(t)$ (sympy), with no floating point except the root-modulus
figures in \S4. Scripts and logs: \texttt{scratch/green-1007c1/} (this session:
\texttt{c1\_resolve.log}, \texttt{c1\_newcontrols.log}, \texttt{newlemmas.log},
\texttt{task2\_census.log}, \texttt{task2\_thmD.log}, \texttt{task2\_oddb.log},
\texttt{task3\_bridge2.log}, \texttt{task3\_bridge\_final.log}, \texttt{thmC8.log},
\texttt{verify\_all\_1007.log}, \texttt{raising\_xcheck.log}) and
\texttt{work-in-progress/code-1006c2-green/} (2026-10-06 c2: \texttt{lemma1.log},
\texttt{vacuity.log}, \texttt{controls.log}, \texttt{xcheck\_AB.log}). The Hall--Littlewood cache
for $n\le 9$ was built on 2026-10-06 and reused unchanged.

Theorem~\ref{thm:C}'s verification log, absent on 2026-10-06, was regenerated this session
(\texttt{thmC8.log}, $78$ checks at $n\le 8$, and $110$ at $n\le 9$ in
\texttt{verify\_all\_1007.log}).

\end{document}
